\section{The closure theorem}\label{sec:final}

We collect the results of the paper in one theorem. Its first clause is
\Cref{thm:main}; the second and third give the base of that theorem in its
sharpest form; the fourth and fifth say what closing the algebraic locus
requires and which constructions cannot supply it; and the last places the
Weil classes inside the Hodge conjecture and states what \Cref{thm:main}
becomes when it is carried to every family. The theorems we quote from the
literature are used as theorems: the semiregularity theorem of Buchweitz and
Flenner \cite{BF03}, in the form for perfect complexes that Pridham proves
\cite{Pri24}; the theorem of Cattani, Deligne and Kaplan \cite{CDK95}; the
theorem of Baily and Borel \cite{BB66}; the theorem of Pila, Shankar and
Tsimerman on the Andr\'e--Oort conjecture \cite{PST}; the description of the
Hodge rings of abelian varieties of dimension at most five by Moonen and Zarhin
\cite{MZ99}; Markman's theorems on abelian varieties of Weil type
\cite{Mar25a,Mar25b}; the theorem of Bando and Siu on stable reflexive
sheaves \cite{BS94}, Kuranishi's theorem in the form of Goldman and Millson
\cite{Kur65,GM88}, Grivaux's Chern classes of coherent sheaves \cite{Gri10}
and Voisin's theorem on Weil tori \cite{Voi02b}; and the theorems of Deligne,
Andr\'e and Milne recorded in \Cref{ssec:bullseye}
\cite{Del71,And96b,Mil20}.

\begin{theorem}[Closure]\label{thm:final}
Let $F$ be a CM field of degree $2m$, let $n\ge2$, and let $\delta$ be a
discriminant. Write $K=\QQ(\sqrt{-d})$ when $m=1$.
\begin{enumerate}[label=\textup{(\roman*)},leftmargin=2.8em,itemsep=4pt]
\item \emph{The algebraic locus.} The algebraic locus $\Sigma_{F}$ of the Weil
  classes contains an explicit CM point at which they are polynomials with
  rational coefficients in divisor classes, it is dense, it is a countable
  union of closed algebraic strata, and it is either the whole domain or
  meagre. It is closed if and only if the Hodge conjecture holds for the very
  general member of the family, and then the Weil classes are algebraic on
  every member (\Cref{thm:main}).
\item \emph{One class on one connected domain.} For an abelian $2n$-fold $A$ of
  $(K,-1,n)$-Weil type with $\Hg(A)=\SU(V,H)$ one has
  $\dim_{\QQ}\Hdg^{k}(A)=1$ for $k\ne n$ and $3$ for $k=n$, and the Hodge
  conjecture for $A$ in every degree is the algebraicity of the single class
  $\omega_{1}$ of \eqref{eq:omegagens}, written out with $2^{2n-1}$ integer
  coefficients. For the family of the triple $(K,n,\delta)$ the conjecture is
  equivalent to the algebraicity of the absolutely Hodge class $\omega_{s}$ at
  every point $s$ of the connected period domain $\cD_{n,n}$ of dimension
  $n^{2}$ (\Cref{thm:hdgcount}, \Cref{cor:whatisleft}, \Cref{thm:residual},
  \Cref{cor:exactform}).
\item \emph{A base point in every family.} Every family of every CM field
  contains a member isogenous to a product of CM abelian $m$-folds at which
  the Weil classes are polynomials in divisor classes
  (\Cref{thm:cmbasepoint}). For an imaginary quadratic field every family also
  contains a member with quaternionic multiplication extending $K$, at which
  the Weil classes are the real and imaginary parts of the $n$-th power of one
  explicit $(1,1)$ class, and every product of $n$ abelian surfaces of Weil
  type is a member whose Weil classes are an explicit integral polynomial in
  $2n$ divisor classes (\Cref{thm:quatweil}, \Cref{thm:everyfamily},
  \Cref{thm:everyn}).
\item \emph{Closing the locus.} If at every point of the orbit
  $G_{F}(\QQ)\cdot s_{0}$ the Weil class is represented by a cycle whose
  Hilbert data comes from one finite set, then $\Sigma_{F}=\cD_{F}$; this
  hypothesis is equivalent to its conclusion, and the transported cycle does
  not satisfy it whenever a component of the base cycle has finite stabiliser
  (\Cref{thm:degreebound}, \Cref{thm:p1equivalent}, \Cref{thm:p1forces}).
  The bound is needed only along one sequence of CM points that leaves every
  proper special subvariety, and such a sequence exists inside the orbit
  (\Cref{thm:escape}). The locus of members with smaller Hodge group, which
  contains every member at which this paper exhibits algebraic Weil
  classes, contains the base point, is dense, stable under
  $G_{F}(\QQ)$ and a countable union of closed algebraic strata, and is
  meagre, so closing the locus requires the algebraicity of the Weil classes
  at a member with full Hodge group (\Cref{prop:speciallocus}, \Cref{cor:speciallocus}). If
  at a base point some perfect complex whose Chern character is a nonzero
  multiple of a rational Weil class plus a polynomial in the divisor classes
  of the family has injective semiregularity map, or satisfies the flatness
  condition of \Cref{prop:p2flat}, then $\Sigma_{F}=\cD_{F}$
  (\Cref{thm:semiregclosure}, \Cref{thm:cmpropagation}). When the Chern
  character is exactly a multiple of a Weil class, no complex at any member of
  any family, for any CM field, has injective semiregularity map
  (\Cref{thm:p2false}, \Cref{thm:cmpurefalse}). In the polarised form, for an
  imaginary quadratic field and $n\ge3$, every object has
  $\dim\Ext^{2}(E,E)\ge(4+\rho)n^{2}-2n$, with $\rho\le3$ the rank of the
  Hankel matrix of the coefficients, and an object attaining the bound meets
  the criterion (\Cref{thm:p2primenumber}); at a quartic
  CM field and $n=2$ the number takes fifteen values, an object meeting the
  criterion has $\dim\Ext^{2}(E,E)\ge88$, and the characters of rank $88$
  whose Hodge locus leaves the polarised family are exactly those of
  $\theta^{4}$ shape, which the Weil tori exclude (\Cref{thm:quarticrank},
  \Cref{cor:quarticleast}, \Cref{prop:quarticlocus},
  \Cref{cor:quarticlocus}). The statement is needed only for the split
  families of the CM fields of degree at least four, which carry the Weil
  classes of every CM field by descent and scalar extension
  (\Cref{prop:descent}), and for one family it is equivalent to the Lefschetz
  standard conjecture for the total space of the family over a curve through a
  base point (\Cref{cor:weilequiv}).
\item \emph{Constructions that do not reach a Weil class.} Each of the
  following fails to produce a Weil class on a general member, in the sense
  and under the hypotheses of the theorem cited, which \Cref{sec:scope}
  delimits:
  \begin{enumerate}[label=\textup{(\alph*)},nosep,leftmargin=2.2em]
  \item an object restricted from a projective space to the dual and
    transported by the Poincar\'e kernel (\Cref{thm:divisor});
  \item equivariant for the field with the norm character
    (\Cref{thm:character});
  \item an object whose Chern character lies on the Weil line and on which
    $H^{2}(\cO_{A})$ acts faithfully, whether through semiregularity or
    through the weakened criterion of \cite[Question~11.4]{Mar25b}
    (\Cref{thm:rankzero}, \Cref{cor:kernelfamily}, \Cref{thm:weakfails});
  \item an object with a line bundle summand off the polarisation line, in
    particular a sum of line bundles (\Cref{thm:linebundle},
    \Cref{cor:noline});
  \item an object whose Chern character is effective at a base point
    (\Cref{thm:positivity});
  \item induction through hyperplane sections, in either direction
    (\Cref{thm:hyperplane}, \Cref{thm:returntrip});
  \item Kuga--Satake from a weight two sub-structure of $H^{2}(A,\QQ)$
    (\Cref{prop:normpart}, \Cref{cor:ksboundary});
  \item a divisorial argument at a general member of a family of nontrivial
    discriminant (\Cref{prop:divfails});
  \item a secant object whose support has two smooth branches of codimension
    two meeting transversally at a point, which includes the candidate of
    \cite[Section~12]{Mar25b} (\Cref{thm:candidatefails},
    \Cref{cor:localobstruction}, \Cref{cor:cmdichotomy});
  \item a secant object of rank one whose support is a complete intersection
    of two divisors (\Cref{thm:noci}), or a secant object quasi-isomorphic to
    a two-term complex of sums of powers of the polarisation or of simple
    semi-homogeneous bundles (\Cref{thm:nosplit}, \Cref{thm:noblocks}); on
    the fourfold, a secant sheaf $\cI_{S}(3\Theta)$ with smooth support
    unless $d\le9$, that is $d\in\{1,3,5,7\}$ in \Cref{setup:markman}, and
    one whose ideal sheaf is presented by two bundles
    with classes in the subring generated by line bundles unless
    $d\equiv15,23\pmod{24}$ (\Cref{cor:fourdiscriminants},
    \Cref{cor:latticeburch});
  \item a representative of the Weil $K$-theory class that is a locally free
    sheaf on a smooth germ of codimension at least two of its support, or
    whose support has codimension at least $n$ (\Cref{thm:p2forces},
    \Cref{thm:p2support});
  \item at a quartic CM field, a complex with secant character that is
    locally free on a smooth curve at some point, in particular the explicit
    pair of \cite[Example~11.2.7]{Mar25c} (\Cref{thm:quarticlocal},
    \Cref{cor:markmanquarticfails});
  \item a convolution of fewer than $2n$ line bundles with nondegenerate
    differences (\Cref{thm:fewpieces}); on $E_{0}^{6}$, every convolution
    of the Weil pieces with three multiples of the polarisation, and on
    $E_{0}^{8}$ every one whose pieces have their shifts within six
    consecutive values (\Cref{thm:noconvolution}, \Cref{cor:efoursix}).
  \end{enumerate}
\item \emph{The Weil classes inside the conjecture.} The Hodge conjecture is
  equivalent to the conjunction of \textup{(F2)}, the algebraicity of the
  Hodge classes of abelian varieties outside the subring generated by divisor
  classes and Weil classes of quotients, and \textup{(F3$'$)}, the
  algebraicity of every Hodge class modulo the images of Hodge classes of
  abelian varieties under algebraic correspondences; and \textup{(F2)} alone
  is equivalent to the Hodge conjecture for abelian varieties
  (\Cref{prop:f2isab}, \Cref{prop:f3prime}). The statement \textup{(F3$'$)}
  holds on every variety whose cohomology is reached from abelian varieties by
  correspondences, on every variety of dimension at most three, on every
  fourfold and fivefold whose zero-cycles are supported in dimension at most
  three or that is dominated by a variety of the first kind, on every smooth
  ample divisor of odd dimension in a variety of the first kind, and on every
  smooth Delsarte hypersurface of every dimension and every cyclic cover
  branched along one, and on every complete intersection of diagonal
  hypersurfaces of Vandermonde type, where the Hodge conjecture follows for
  two diagonal hypersurfaces of degree $3$, $4$ or $6$ in $\PP^{6}$ and for
  the very general member of degree $2$ in every dimension, of degree $3$ or
  $4$ up to dimension seven and of degree $6$ up to dimension five
  (\Cref{prop:f3prime}, \Cref{prop:aclosure},
  \Cref{prop:f3primesmall}, \Cref{prop:f3primefibration},
  \Cref{prop:f3primedominant}, \Cref{prop:f3primeample}, \Cref{cor:delsarteall},
  \Cref{thm:vandermonde}, \Cref{cor:twodiagonal},
  \Cref{thm:vgvandermonde}). The first instance of
  \textup{(F2)}, the two exceptional classes on the square of a Mumford
  fourfold, is algebraic at every CM point of the family, is algebraic on
  every member once one Kuga--Satake class is algebraic, and is algebraic on
  every member exactly when the Lefschetz standard conjecture holds for one
  ninefold (\Cref{thm:mumfordcm}, \Cref{thm:mumfordks},
  \Cref{cor:mumfordequiv}). Over the rule set formed by the implications of
  the paper and of the literature, together with \Cref{prop:f2isab}, the
  minimal sets of statements from which the conjecture follows are
  $\{\mathrm{F3}\}$, $\{\mathrm{L},\mathrm{M}\}$,
  $\{\mathrm{L},\mathrm{F3}'\}$, $\{\mathrm{V},\mathrm{F3}'\}$ and
  $\{\mathrm{F2},\mathrm{F3}'\}$, the first two being each equivalent to it
  (\Cref{thm:closure}, \Cref{rem:f2rule}). Carried to every family,
  \Cref{thm:main} gives exactly these two statements back:
  \textup{(F2)} is equivalent to the statement that in every polarised
  abelian scheme with a member of CM type the algebraic locus of a flat
  Hodge class that is algebraic at such a member has an interior point
  (\Cref{thm:f2propagation}); \textup{(F3$'$)} is equivalent to the absence
  of Hodge classes off the part of the cohomology reached from abelian
  varieties by algebraic correspondences (\Cref{prop:abpart}); and on the
  squares of K3 surfaces with real multiplication, the first varieties on
  which \textup{(F3$'$)} is open, the algebraic locus of the real
  multiplication is dense, contains every CM point and is either the whole
  family or meagre (\Cref{thm:rmlocus}).\footnote{These are the statements
  \textup{(L)}, \textup{(M)} and \textup{(V)} of \Cref{def:lmv}: $B(X)$ of
  \Cref{def:Bconj} for every smooth projective $X$; every Hodge class
  motivated in the sense of Andr\'e \cite[\S2.1]{And96b}; and, for every
  smooth projective morphism $f$ to a smooth connected base, a flat section of
  $R^{2p}f_{*}\QQ$ that is algebraic at one point is algebraic at every
  point. The statement (F3) is the Hodge conjecture for the varieties that
  are not abelian, which is equivalent to the whole conjecture
  (\Cref{prop:f3ishc}).}
  With the theorems of the preprints of \Cref{sec:preprints} as
  hypotheses, the two exceptional classes are algebraic on every member,
  the Weil classes are algebraic on every abelian sixfold of Weil type, the
  real multiplication is algebraic on every K3 surface, and neither
  \textup{(F2)} nor \textup{(F3$'$)} follows (\Cref{cor:mumfordall},
  \Cref{thm:weilsix}, \Cref{cor:k3all}, \Cref{rem:remains}).
\end{enumerate}
\end{theorem}

\begin{proof}
Each clause is the theorem or theorems cited with it, and no clause is used
in the proof of another. For (i), \Cref{thm:main} is proved at the end of
\Cref{ssec:cmpropagation}. For (ii), \Cref{thm:hdgcount} counts the Hodge
classes, \Cref{cor:whatisleft} reduces every degree to $\omega_{1}$,
\Cref{thm:explicitgens}(i) shows that $\omega_{1}$ has integer coefficients
and exactly $2^{2n-1}$ nonzero terms, \Cref{thm:residual}(a) and (b) are the
two implications for the family, and \Cref{prop:weildomain} gives the
connectedness and the dimension. For (iii), \Cref{thm:cmbasepoint} gives the
CM point, \Cref{thm:quatweil}(v) the cycle at a quaternionic point,
\Cref{thm:everyfamily} places such a point in every family by the
discriminant computation of \Cref{thm:quatweil}(iv) and \Cref{lem:products},
and \Cref{thm:everyn} treats products.

For (iv), the bounded criterion is \Cref{thm:degreebound} with
\Cref{thm:p1equivalent} and \Cref{thm:p1forces}, and its form for a CM field
is \Cref{thm:cmpropagation}(iv); the escaping sequence is \Cref{thm:escape},
and the special locus is \Cref{prop:speciallocus} with
\Cref{cor:speciallocus}. The semiregularity criterion is
\Cref{thm:semiregclosure} with \Cref{rem:p2prime}, carried to every CM field
by \Cref{thm:cmpropagation}(v); the flatness form is \Cref{prop:p2flat}. The
pure form is excluded by \Cref{thm:p2false}, which rests on
\Cref{lem:weiltorushodge}, \Cref{lem:weiltorussubsets} and
\Cref{prop:weiltorussheaves}, and for every CM field by
\Cref{thm:cmpurefalse}. The number of the polarised form is
\Cref{thm:p2primenumber}(i),
(iv) and (v); the quartic rank and the Hodge loci are
\Cref{thm:quarticrank}, \Cref{cor:quarticleast}, \Cref{prop:quarticlocus}
and \Cref{cor:quarticlocus}. The reduction to the split families is
\Cref{prop:descent}, and the equivalence with the Lefschetz standard
conjecture is \Cref{cor:weilequiv}. For (v), each item is the statement of
the theorem cited, under its own hypotheses, which \Cref{sec:scope} delimits
item by item.

For (vi), \Cref{prop:f3prime}(i) gives the Hodge conjecture from
\textup{(F3$'$)} and the conjecture for abelian varieties, and
\Cref{prop:f2isab} gives the latter from \textup{(F2)}; the converse
implications hold because \textup{(F2)} is a case of the conjecture and
\textup{(F3$'$)} is weaker than it. The cases of \textup{(F3$'$)} are
\Cref{prop:f3prime}(iii) and (iv), \Cref{prop:aclosure},
\Cref{prop:f3primesmall}, \Cref{prop:f3primefibration},
\Cref{prop:f3primedominant}, \Cref{cor:delsarteall},
\Cref{thm:vandermonde}, \Cref{cor:twodiagonal} and \Cref{thm:vgvandermonde}.
The statements on
the Mumford classes are \Cref{thm:mumfordcm}, \Cref{thm:mumfordks}(ii) and
\Cref{cor:mumfordequiv}. The statements on the families are
\Cref{thm:f2propagation}, \Cref{prop:abpart} and \Cref{thm:rmlocus}. The
conditional statements are those of \Cref{sec:preprints} cited with them,
and no other clause uses them. The
minimal sets are \Cref{thm:closure}(ii) with
the rule of \Cref{prop:f2isab} adjoined, as computed in \Cref{rem:f2rule},
and \Cref{thm:closure}(vi) shows that $\{\mathrm{F3}\}$ and
$\{\mathrm{L},\mathrm{M}\}$ are each equivalent to the conjecture.
\end{proof}

The theorem arranges the paper around \Cref{thm:main}. Clause (ii) is the
problem in its smallest form, one explicitly written class on one connected
domain, and clause (iii) shows that the base case is never the difficulty:
every family has a base point, for every field, every dimension and every
discriminant, with a cycle written down rather than an existence statement.
Clause (iv) reduces the closedness of the locus to one object at one point
and computes the number that object must attain, and clause (v) shows that
the constructions that fail do so for identifiable structural reasons: a
divisor-generated Chern character cannot leave the divisor subring, a norm
eigenclass sits in the wrong summand of the character grading, a Chern
character on the Weil line has an annihilator exactly as large as the family,
a support with two transversal branches of codimension two through a point
carries a local obstruction that the compensating scalar cannot meet, and a curve along which an object is locally
free is tangent to at most one of the two eigenplanes of the real
multiplication. Clause (vi) shows where the Weil classes stand in the whole
conjecture: the conjecture for abelian varieties is carried by the classes
beyond the Weil lines, and the rest is algebraicity modulo abelian varieties,
which holds on every variety whose cohomology comes from abelian varieties,
on the Delsarte hypersurfaces of every dimension and on the diagonal
complete intersections of Vandermonde type (\Cref{fig:twohalves}).
\Cref{sec:families} shows that the structure \Cref{thm:main} finds in the
Weil families, a CM seed, density and the dichotomy, is present in every
family that carries (F2) or \textup{(F3$'$)}, and that in each of them the
one missing property is the interior point that closes the locus
(\Cref{thm:twostand}).

\begin{figure}[tbp]
\centering
  \includegraphics[width=\textwidth]{fig_twohalves}
\caption{The two halves of the conjecture in \Cref{thm:final}(vi). A Hodge
class $\gamma$ on a smooth projective variety $X$ is the image
$\Gamma_{*}\beta$ of a Hodge class $\beta$ on an abelian variety $B$ under an
algebraic correspondence $\Gamma$ on $B\times X$, which is \textup{(F3$'$)};
the class $\beta$ is algebraic, which is \textup{(F2)}, equivalently the
conjecture for abelian varieties (\Cref{prop:f2isab}). The raised plate on
the lower sheet stands for the families of Weil type, on which
\Cref{thm:main} makes the algebraic locus of the Weil class $\omega$ dense;
the dashed part of each correspondence is the part behind the upper sheet.}
\label{fig:twohalves}
\end{figure}
